\documentclass[12pt]{article}
\usepackage[utf8]{inputenc}
\usepackage[T1]{fontenc}
\usepackage{hyperref}
\title{Nubaca Language Resources for K-12 Curriculum Development}
\author{Compiled by Claude Sonnet 5.5 (claude-sonnet-5-5) for Andr\'as Kornai}
\date{}
\begin{document}
\maketitle

\section{Introduction}
Nubaca (ISO 639-3 baf; Glottocode nuba1241) is a Bantu language of central Cameroon, also called Baca, Bango or Bongo, spoken around Bongo village in the Bokito subdivision. Speaker counts in the sources found range from about 4,500 to about 8,400. Wikitongues classes it as threatened and unwritten. Dialects listed are Central, Kelende and Nibieg. Very little educational material was found in search.

\section{Primary Educational Resources}
No primers, textbooks or school programmes for Nubaca were found in search results. Joshua Project notes that Scripture is available in French and Ewondo but not in Nubaca.

\section{Linguistic and Reference Materials}
\begin{itemize}
\item Wikipedia, ``Baca language'': overview of the language. \url{https://en.wikipedia.org/wiki/Baca_language}
\item Sebineni, Alphonsine Flore (2008), a Universit\'e de Yaound\'e I master's thesis on the noun class system, as listed in Glottolog. \url{https://glottolog.org/resource/reference/id/552065}
\item ASJP word list entry for Nubaca. \url{https://asjp.clld.org/languages/NUBACA}
\item Glottolog languoid page. \url{https://glottolog.org/resource/languoid/id/nuba1241.bigmap.html}
\end{itemize}

\section{Religious and Community Resources}
\begin{itemize}
\item Joshua Project language page (baf). \url{https://joshuaproject.net/index.php/languages/baf}
\item Joshua Project, Bango in Cameroon. \url{https://www.joshuaproject.net/people_groups/10652/CM}
\end{itemize}

\section{Digital and Audio Resources}
\begin{itemize}
\item Global Recordings Network page for Nubaca (baf). \url{https://globalrecordings.net/en/language/baf}
\item Wikitongues page for baf. \url{https://wikitongues.org/languages/baf/}
\end{itemize}

\section{Curriculum Development Context}
The language is reported as unwritten and not known to be taught in schools. A curriculum would first require an orthography, a lexicon and basic descriptive work, building on the noun class thesis above, in partnership with the community. Instruction would in practice remain in French and Ewondo-influenced settings until such materials exist.

\end{document}